
\section{Introduction}
\label{sect:intro}


\begin{figure}[t]
    \centering
    \includegraphics[width=\linewidth]{figure/1-introduction/teaser.pdf}
    \vspace{-18pt}
    \caption{\system achieves higher throughput when running Llama models on L40S compared with TensorRT-LLM on A100, effectively saves the dollar cost for LLM serving by \textbf{3\x} through system-algorithm codesign. See Table~\ref{tab:abs_main_speed} for absolute throughput numbers and precision choices in TensorRT-LLM.} %
    \label{fig:intro:teaser}
    \vspace{-18pt}
\end{figure}


Large language models (LLMs) have demonstrated remarkable capability across a broad spectrum of tasks, exerting a profound influence on our daily lives. 
However, the colossal size of LLMs makes their deployment extremely challenging, necessitating the adoption of quantization techniques for efficient inference. State-of-the-art integer quantization algorithms can be divided into three categories: 8-bit weight and 8-bit activation (\textbf{\texttt{W8A8}}), 4-bit weight and 16-bit activation (\textbf{\texttt{W4A16}}), 4-bit weight 4-bit activation (\textbf{\texttt{W4A4}}) quantization. The former two methods are considered nearly lossless in terms of accuracy. In contrast, \textbf{\texttt{W4A4}} quantization introduces a notable accuracy degradation, although it is anticipated to offer superior throughput in return by mapping its computations onto high-throughput 4-bit tensor cores. Unfortunately, this anticipated performance boost has not been consistently observed across current GPU platforms. %
For instance, the state-of-the-art \textbf{\texttt{W4A4}} serving system, Atom~\cite{zhao2023atom}, exhibits 20-25\% lower performance than its \textbf{\texttt{W4A16}} and \textbf{\texttt{W8A8}} counterpart in TensorRT-LLM when running the Llama-2-7B~\cite{touvron2023llama2} model on A100 GPUs. That said, the research community has yet to find a precision combination superior to \textbf{\texttt{W4A16}} and \textbf{\texttt{W8A8}} for efficient cloud LLM serving. 



In this paper, we reveal a critical observation: current 4-bit integer quantization methods experience significant overhead, ranging from 20\% to 90\%, during the dequantization of weights or partial sums %
on current-generation GPUs. For example, \textbf{\texttt{W4A16}} quantization performs computation on \texttt{FP16} tensor cores while the weights are in \texttt{INT4}, so weight dequantization is required in the GEMM kernel. On the other hand, for \textbf{\texttt{W4A4}} quantization, to achieve reasonable accuracy, \textbf{\texttt{W4A4}} methods must apply \textbf{per-group} quantization to both weights and activation, sharing \texttt{FP16} scaling factors on a sub-channel basis. For example, the state-of-the-art \textbf{\texttt{W4A4}} quantization method, QuaRot~\cite{ashkboos2024quarot}, reports a significant 0.2 perplexity degradation after switching from per-group quantization to per-channel quantization. This per-group quantization design requires an integer to floating-point dequantization for partial sums (since \texttt{INT4} tensor cores produce \texttt{INT32} partial sums), which operates on the slower CUDA cores within the sequential main loop of \textbf{\texttt{W4A4}} GEMM. On data center GPUs like A100, a CUDA core operation is as expensive as \textbf{50} \texttt{INT4} tensor core operations. Therefore, reducing bit precision not necessarily speeds up LLM inference. 

To achieve optimal LLM serving throughput, We introduce \algo (Quattuor-Octō-Quattuor, or 4-8-4 in Latin) algorithm which quantizes LLMs to \textbf{\texttt{W4A8KV4}} precision: 4-bit weights, 8-bit activations and 4-bit KV caches. Additionally, we present \system, which provides efficient system support for \textbf{\texttt{W4A8KV4}} quantization. 

In the \algo algorithm, we introduce \textit{progressive group quantization}. This method first quantizes weights to 8 bits using per-channel \texttt{FP16} scales with a \textit{protective range} of $[-119, 119]$, then quantizes these 8-bit intermediates to 4 bits. This approach ensures that all GEMMs are performed on \texttt{INT8} tensor cores. Additionally, we mitigate accuracy loss from \texttt{KV4} quantization through \textit{SmoothAttention}, which shifts the challenge of activation quantization from keys to queries, the latter of which are not quantized. 

In the \system system, the \textit{protective range} in progressive group quantization enables full \textit{register-level parallelism} during \texttt{INT4} to \texttt{INT8} dequantization, using a \textit{subtraction after multiplication} computation order. Furthermore, we propose \textit{compute-aware weight reordering} to minimize pointer arithmetic overhead on CUDA cores during W4A8 GEMM operations. Additionally, we delay the turning point of the CUDA core roofline and decrease the computational intensity of \texttt{KV4} attention at the same time. This ensures that the attention operator remains within the memory-bound region, where low-bit quantization can effectively enhance throughput.

We evaluate seven widely-used LLMs using \system on A100 and L40S GPUs, and compare their maximum achievable throughput against state-of-the-art systems, including TensorRT-LLM (in \texttt{FP16}, \textbf{\texttt{W8A8}}, and \textbf{\texttt{W4A16}} configurations), Atom~\cite{zhao2023atom} (in \textbf{\texttt{W4A4}}), and QuaRot~\cite{ashkboos2024quarot} (in \textbf{\texttt{W4A4}}). On A100 GPUs, \system achieves \textbf{1.2-2.4\x} higher throughput over the best-performing configuration of TensorRT-LLM, and \textbf{2.5-2.9\x} higher throughput compared to Atom and QuaRot. On L40S GPUs, \system records an even more significant \textbf{1.5-3.5\x} throughput improvement over TensorRT-LLM. Notably, we manage to accommodate the same batch size on the L40S while consistently achieving higher serving throughput than TensorRT-LLM on A100 which is 3$\times$ more expensive for six of the eight models tested.




